\section{Typed ground action specifications}\label{sec:setup}

\subsection{Native diagrams and the action symbol}

All algebras are nonempty and all native operation signatures are finite and finitary.  Elements are named by designated constants; native operation symbols are assumed to have positive arity.  Let
\[
A=(|A|,\Sigma_A),\qquad B=(|B|,\Sigma_B)
\]
be finite algebras.  The two-sorted signature has sorts $\A$ and $\B$, native operations of $\Sigma_A$ on $\A$, native operations of $\Sigma_B$ on $\B$, designated constants $c_a:\A$ for $a\in A$ and $c_b:\B$ for $b\in B$, and
\[
\alpha:\B\times\A\to\A.
\]
We usually suppress the constant symbols and write $a,b$ for their names.

The complete native ground diagrams are
\[
D_A=\{f(a_1,\ldots,a_k)=f^A(a_1,\ldots,a_k)\},
\]
with one row for every $f\in\Sigma_A$ and named tuple, and similarly $D_B$.  The base action theory contains the ground carrier-compatibility instances
\begin{equation}\label{eq:compat}
\alpha(b,f(a_1,\ldots,a_k))
=
f(\alpha(b,a_1),\ldots,\alpha(b,a_k))
\end{equation}
for named $b\in B$ and named $a_i\in A$.

An incomplete action table is an arbitrary finite set of row occurrences
\begin{equation}\label{eq:rows}
\Gamma\subseteq
\{\alpha(b,a)=c:b\in B,\ a,c\in A\}.
\end{equation}
Occurrences, rather than a partial function, are allowed: two supplied equations may have the same input and different outputs.  A complete table is the special case in which one output is supplied for every $(b,a)$.

Sections~\ref{sec:rows}--\ref{sec:repair} study the \emph{pure row fragment}
\[
E=D_A\cup D_B\cup\mathrm{Comp}_A\cup\Gamma,
\]
where $\mathrm{Comp}_A$ denotes the family \eqref{eq:compat}.  Ground action laws such as unit, composition, and inverse are added later.

These are equations between closed terms, not universally quantified
identities with unnamed inputs. Models range over all algebras of the
specified two-sorted signature. In particular, an identity that happens to
hold in $A$ is not automatically required on the whole carrier sort of a
model. The operator sort of the initial algebra is exactly its named copy
of $B$: every closed operator term reduces by $D_B$, and the model
constructions below separate its names.


\subsection{Initial algebra, carrier quotient, and external values}

Let $\mathbb I_E$ be the initial two-sorted algebra of the equational theory $E$.  The map from carrier names into the $\A$-sort of $\mathbb I_E$ need not be injective.  Define the induced carrier congruence
\begin{equation}\label{eq:thetaE}
\theta_E^A
=
\{(a,a')\in A^2:E\vdash a=a'\}.
\end{equation}
Since provable equality is a congruence, $\theta_E^A\in\Con(A)$.  Put
\[
Q_E=A/\theta_E^A.
\]
The named carrier embeds canonically as $Q_E$ in the initial $\A$-sort.

\begin{definition}
For a named $b\in B$ and $a\in A$, the class of $\alpha(b,a)$ is \emph{carrier-valued} if it contains a carrier name.  Otherwise it is \emph{external}.  The named ground action is \emph{closed} if every $\alpha(b,a)$ is carrier-valued.
\end{definition}

The three central diagnostics are now formal.

\begin{definition}[Protection, determinacy, realizability]
The presentation is \emph{carrier-protecting} when $\theta_E^A=\Delta_A$.  For a fixed row $b$, its \emph{forced domain} in $Q_E$ is
\[
D_b=
\{[a]\in Q_E:\exists c\in A\text{ with }E\vdash\alpha(b,a)=c\}.
\]
A \emph{carrier-valued completion} on a quotient $A/\theta$ is a family of endomorphisms
\[
\tau_b\in\End(A/\theta)\qquad(b\in B)
\]
satisfying all projected supplied rows.  Additional action laws, when present, impose the corresponding equations on the family $(\tau_b)$.
\end{definition}

The endomorphism requirement is exactly carrier compatibility: a total row map $\tau_b:Q\to Q$ satisfies all instances of \eqref{eq:compat} precisely when it preserves every native carrier operation.

\begin{definition}[Repair set]
For a pure-row presentation, define
\begin{equation}\label{eq:repairset}
\mathcal R(E)=
\{\theta\in\Con(A):\text{the projected rows admit a carrier-valued completion on }A/\theta\}.
\end{equation}
Elements of $\mathcal R(E)$ are called \emph{repairs}.
\end{definition}

The total congruence is always a repair, so $\mathcal R(E)$ is nonempty.


\begin{center}
\begin{tabularx}{\textwidth}{@{}lX@{}}
\toprule
Symbol & Meaning throughout the pure-row sections \\
\midrule
$\theta_E^A$, $Q_E$ & Forced carrier congruence and $A/\theta_E^A$. \\
$S_b$, $h_b$ & Subalgebra generated by supplied inputs, and its induced row homomorphism on a protected quotient. \\
$K_b$, $\kappa_b$ & $\ker h_b$ and its generated congruence $\Cg_Q(K_b)$ on $Q$. \\
$D_b$ & Domain of the carrier-valued part of the initial row. \\
$\theta^{[r]}$ & Carrier congruence after $r$ feedback rounds. \\
$D$, $\equiv_D$ & Term-depth horizon and equality visible at that horizon. \\
$\mathcal I$ & Named carrier and named action interface. \\
$\mathcal R(E)$ & Quotients admitting a carrier-valued completion. \\
$\kappa_E^{\mathrm{NJ}}$ & Intersection of all repair congruences. \\
\bottomrule
\end{tabularx}
\end{center}
